\documentclass[11pt,a4paper]{article}
\newcommand{\fichetype}{Exercices}
% Thème A4 commun, librement modifiable.
\usepackage[T1]{fontenc}
\usepackage[utf8]{inputenc}
\usepackage[french]{babel}
\usepackage{lmodern,amsmath,amssymb,mathtools,tikz,pgfplots}
\usepackage[margin=2.2cm,headheight=15pt]{geometry}
\usepackage{xcolor,fancyhdr,enumitem,hyperref,needspace}
\usetikzlibrary{arrows.meta,positioning,calc}
\pgfplotsset{compat=1.18,width=11cm,height=6cm}
\definecolor{Ink}{HTML}{202124}
\definecolor{Blue}{HTML}{356B9E}
\definecolor{Muted}{HTML}{737579}
\color{Ink}
\renewcommand{\familydefault}{\sfdefault}
\setlength{\parindent}{0pt}
\setlength{\parskip}{6pt}
\setlist{itemsep=3pt,topsep=4pt}
\pagestyle{fancy}\fancyhf{}
\fancyhead[L]{\small\color{Muted}MATHS LYCÉE}
\fancyhead[R]{\small\color{Muted}2026--2027}
\fancyfoot[L]{\small\color{Muted}\fichetype}
\fancyfoot[R]{\small\thepage}
\renewcommand{\headrulewidth}{0.3pt}
\hypersetup{colorlinks=true,urlcolor=Blue,linkcolor=Blue}
\newcommand{\R}{\mathbb{R}}
\newcommand{\N}{\mathbb{N}}
\newcommand{\E}{\mathbb{E}}
\newcommand{\V}{\mathrm{V}}
\newcommand{\exercice}[3]{\Needspace{6\baselineskip}\vspace{0.5cm}{\color{Blue}\small\textbf{EXERCICE #1}\hfill #2}\par{\large\bfseries #3}\par}

\begin{document}
{\small\color{Blue}PREMIÈRE SPÉCIALITÉ}\par{\LARGE\bfseries Suites arithmétiques et géométriques}\par{\large Exercices}\par\vspace{0.2cm}\hrule\vspace{0.4cm}
Les exercices progressent de l’application directe vers la recherche. Détailler les calculs et justifier les réponses.
\par\noindent\begin{minipage}{\linewidth}\exercice{1}{Application}{Une suite arithmétique}
$u_0=12$ et $r=-3$. Calculer $u_7$, écrire $u_n$ et déterminer le premier rang tel que $u_n<0$.
\end{minipage}\par
\par\noindent\begin{minipage}{\linewidth}\exercice{2}{Entraînement}{Retrouver la raison}
Une suite arithmétique vérifie $u_3=11$ et $u_8=31$. Trouver $r$, $u_0$ et $u_{20}$.
\end{minipage}\par
\par\noindent\begin{minipage}{\linewidth}\exercice{3}{Application}{Sommes arithmétiques}
Calculer $1+\cdots+80$, puis $7+11+15+\cdots+63$.
\end{minipage}\par
\par\noindent\begin{minipage}{\linewidth}\exercice{4}{Application}{Une suite géométrique}
$v_1=6$, $q=1/2$. Calculer $v_5$, puis $v_1+\cdots+v_5$.
\end{minipage}\par
\par\noindent\begin{minipage}{\linewidth}\exercice{5}{Entraînement}{Variations comparées}
Étudier les variations de $a_n=4(0,7)^n$, $b_n=-3\times2^n$ et $c_n=5(-1/2)^n$.
\end{minipage}\par
\par\noindent\begin{minipage}{\linewidth}\exercice{6}{Entraînement}{Placement et versement}
Un capital fictif de $1500$ euros augmente de $4\,\%$ par an sans versement. Donner sa valeur après trois ans et le taux total d’évolution. Comparer à trois hausses additionnées.
\end{minipage}\par
\par\noindent\begin{minipage}{\linewidth}\exercice{7}{Entraînement}{Algorithme de cumul}
On initialise $v=2$, $S=0$. On répète six fois : $S\leftarrow S+v$, puis $v\leftarrow3v$. Quelle somme calcule-t-on ? La calculer sans boucle.
\end{minipage}\par
\par\noindent\begin{minipage}{\linewidth}\exercice{8}{Approfondissement}{Remboursement : changement de suite}
Une dette fictive suit $D_{n+1}=1,1D_n-100$, avec $D_0=500$. Poser $v_n=D_n-1000$. Montrer que $v$ est géométrique et exprimer $D_n$, tant que le modèle décrit une dette.
\end{minipage}\par
\end{document}
